\section{总结}
确定性条件下，储能的价值来自将有限电量移至更需要的时段；在给定效率、日循环和余电处置条件下，问题一最低日费为35126.95元。存在预测误差时，购电问题进一步成为提前支付常规电费与事后支付高价紧急电费之间的权衡。固定价格和波动价格下，所构造日内策略的正式期费用分别下降4.26\%和4.30\%。

决定这种改善的因素应分开判断：完整调整策略的收益大于单独新增预报的收益，18点新预报未显示可辨认的经济增益；增购不能撤回时，收益仍存在但缩小。较大风险分位并非总是更省，逐月重选参数和所测试的价格感知反馈也未改善累计费用。因此，后续应用应同时核对信息、结算和库存边界，而不只比较一个总费用数字。
\label{end:body}
